\chapter{Conclusions}\label{ch:conclusions}

This thesis characterised three causal linear-time recurrent
backbones (RWKV-6, Mamba-2, Linear Attention) and three
mechanism-level extensions (Multi-Scale Depthwise Convolution,
Damped Harmonic Oscillator, Chunked Value Decorrelation) on a
cross-architecture transfer matrix at 7M and 14M parameter
scales. The matrix supports a predictive criterion that we call
mechanism-prior alignment: a mechanism's gain on a task is
proportional to the residual structural deficit the architecture
leaves uncovered along the axis the mechanism targets, and is
null when the mechanism's axis is not exercised by the task.

\section{Summary of findings}\label{sec:conclusions:summary}

The matrix populates the criterion in six structural signatures
across modes, scales, probes, and composition settings
(Section~\ref{sec:experiments:synthesis}): deficit-proportional
ordering on MSDC, structural BREAK on DHO on Linear Attention
together with a near-null on Mamba-2 under the LibriSpeech
distribution, asymmetric transfer on CVD, the
LION-LIT~$\times$~CVD controlled falsifier isolating decay as
the bidirectional-mode prerequisite, the
Mamba-2~$\times$~DHO cross-distribution refinement to alignment
between mechanism and task-as-exercised-by-data, and the axis-2
isolation on the Multi-Query Associative Recall length sweep
where mechanism-augmented linear-time backbones solve the task
at $T = 1024$ while the causal Transformer fails. The matrix
ceiling on LibriSpeech is the 14M Mamba-2 CVD~$\times$~MSDC
composition cell at test character error rate $0.0618$; the
strongest single-mechanism gain is the LA LION-LIT~$\times$~DHO
BREAK at a $65\%$ relative reduction. The chapter does not
lead on these absolute or relative numbers; the artefact is the
matrix and the predictive criterion the matrix supports.

\section{Limitations}\label{sec:conclusions:limitations}

Three limitations bracket the empirical scope. First, the matrix
is single-seed for every reported cell; the across-matrix
pattern is read off approximately seventy statistical tests
oriented in a single coherent direction, but per-cell standard
errors are not estimated. Multi-seed validation of the
BREAK-band cells (LA causal DHO, LA LION-LIT DHO) would sharpen
the within-cell claim without changing the across-matrix
pattern. Second, the matrix scales to 14M parameters; the
conditional 30M scale was dropped from scope since no ranking
inversion appears between 7M and 14M, but the deficit-proportional
ordering is not asserted at scales beyond 14M. Third, the
acoustic probes are English-only; the Common Voice
cross-distribution validation exercises within-language
acoustic variability rather than cross-language transfer.

\section{Scope and future work}\label{sec:conclusions:future-work}

The criterion as stated holds within the diagonal-class function
family of Section~\ref{sec:bounds}. The matrix does not test it
on diagonal-plus-low-rank architectures, and the criterion is
not asserted on that class. The natural extension of the
framework is to repeat the cross-architecture transfer matrix
on a diagonal-plus-low-rank backbone class such as
DeltaNet~\cite{schlag2021deltanet, yang2024deltanet} or
RWKV-7~\cite{peng2025rwkv7}, with mechanisms targeting axis~3
(non-Abelian state-tracking). The formal complexity-class
boundary recalled in Section~\ref{sec:bounds} brackets exactly
where the diagonal class cannot reach and where the
diagonal-plus-low-rank class is the principal route across.

A second extension direction is the cross-modality replication of
the matrix on vision and language probes. The thesis instantiated
MSDC on speech as one node in a five-paper cross-domain pattern
(Section~\ref{sec:related:cross-modality}); the
cross-architecture transfer matrix is portable to those
modalities with appropriate substitution of the axis-1 and
axis-2 probes.

The framework remains agnostic to the specific evaluation domain
and to the specific mechanism instantiation. What it asserts is
the predictive structure: a mechanism's gain orders with the
residual structural deficit along the axis it targets, modulated
by what the task-as-exercised-by-data exercises. The empirical
matrix populates the criterion in three productive signatures
across the diagonal-class linear-time family and one controlled
falsifier; further work extends the criterion across the
diagonal-plus-low-rank boundary and across modalities other
than speech.
